\chapter{讨论}\label{chap:discussion}

\section{核函数选择的物理解释}\label{sec:physics-interp}

六种核函数的 $\xi(s)$ 性能排序揭示了清晰的物理规律：
$\mathrm{Mat\acute{e}rn}(5/2)\approx\mathrm{Mat\acute{e}rn}(3/2)\approx{\rm Combined}
> {\rm RQ} > {\rm RBF} \gg \mathrm{Mat\acute{e}rn}(1/2)$。从偏差-方差权衡角度理解：
RBF 核（$\nu=\infty$，无限可微）先验过于刚性，在有限训练样本（104 组）下导致过度平滑
（欠拟合），增大模型偏差。Mat\'ern 核较弱的光滑性假设（$5/2$ 为二阶可微，$3/2$ 为
一阶可微）提供更灵活的拟合能力，在偏差与方差间取得更好平衡。但过度降低光滑度
（$\nu=1/2$，不可微）则走向另一极端——函数空间过大，模型方差急剧增大，MAE（$0.00821$）
比 RBF 还高 $45\%$。这表明宇宙学关联函数对参数的依赖是光滑的，但不需要无限可微，
二阶可微（$\nu=5/2$）或一阶可微（$\nu=3/2$）即为最佳光滑度。

Mat\'ern$(5/2)$ 与 Mat\'ern$(3/2)$ 在 $\xi(s)$ 上无统计显著差异（$p=0.81$，MAE 仅差
$0.00003$），本文推荐 $\nu=5/2$ 基于以下考量：（a）$\nu=5/2$ 是 GP 文献中最广泛采用的
Mat\'ern 变体\cite{Rasmussen2006}，便于与已有工作对比；（b）Mat\'ern$(5/2)$ 核函数二阶
可微，数值稳定性优于一阶可微的 $\nu=3/2$（后者在核矩阵条件数较大时更易出现数值问题）；
（c）联合约束 FoM 略优（$58.9$ vs $58.3$），表明在 MCMC 参数约束中 $\nu=5/2$ 有微弱优势。
但我们强调，在实际应用中 $\nu=3/2$ 同样是合理选择。RQ 核（MAE${}=0.00539$）介于 RBF 和
Mat\'ern 之间，其多尺度混合特性在当前数据上未带来显著优势。组合核（Mat\'ern $\times$ RBF）
在 $\xi(s)$ 上与 Mat\'ern 表现几乎相同（$p=0.40$），进一步检查核参数（\autoref{tab:best-kernels}）
揭示退化行为：在 $\xi(s)$ 上 RBF 因子 length\_scale ${}=9.8\times 10^{4}$（常数），
退化为近似 Mat\'ern 核。

值得注意的是，归一化策略对核函数排序有显著影响。灵敏度分析表明，在不做归一化
（${\rm norm}={\rm False}$）的情况下，三种核的 $\xi(\mu)$ MAE 几乎相同（RBF: $0.01552$，
Mat\'ern: $0.01551$，Combined: $0.01550$），核函数差异消失。而启用归一化后
（${\rm norm}={\rm True}$），Mat\'ern 核改善至 $0.01532$ 而 RBF 核反而劣化至 $0.01568$，
差异变得显著。这说明归一化通过改变 GP 看到的数据尺度，暴露了不同核函数在建模弱参数依赖
性时的能力差异。因此，数据预处理策略本身也是仿真器设计中需要审慎选择的因素。

\section{\texorpdfstring{$\Omega_c$}{Ωc} 系统偏移与协方差矩阵问题}\label{sec:omega-c-bias}

所有 MCMC 约束的 $\Omega_c$ 中位值在 $0.268$--$0.308$ 之间，而 fiducial 值为 $0.26072$，
存在约 $0.01$--$0.05$ 的偏移。后验预测检验（PPC）显示，$\chi^{2}/{\rm dof}$ 在 fiducial
处 $\xi(s)$ 约为 $0.10$、$\xi(\mu)$ 约为 $0.72$（均小于 $1$），说明协方差矩阵被高估
（$\xi(s)$ 严重高估约 $10$ 倍，$\xi(\mu)$ 高估约 $1.4$ 倍），似然面偏平坦，后验分布主要
由先验形状驱动，$\Omega_c$ 的偏移在统计上并不显著。进一步地，各组约束的 $\Omega_c$ 中位值
（$0.297$--$0.308$）极接近先验均值（$(0.2+0.4)/2=0.300$），RBF $\xi(\mu)$ 的 $\sigma_8$
中位值（$0.856$）也接近先验均值（$(0.54+1.18)/2=0.860$），证实在 $\chi^{2}/{\rm dof}$
约 $0.1$ 的条件下，后验分布本质上由先验形状决定，数据对参数的实际约束极弱。

协方差矩阵高估有两方面原因：
\begin{enumerate}
  \item \textbf{体积不匹配。} 协方差从边长 $400~\mathrm{Mpc}/h$ 的模拟箱估计
    （体积 $6.4\times 10^{7}~(\mathrm{Mpc}/h)^{3}$），而仿真器数据来自 $1000~\mathrm{Mpc}/h$
    箱（体积 $10^{9}~(\mathrm{Mpc}/h)^{3}$），有效体积相差约 $15.6$ 倍。若对协方差矩阵
    施加体积重标定（乘以 $(400/1000)^{3}=0.064$），$\chi^{2}/{\rm dof}$ 将从约 $0.10$ 升至
    约 $1.6$（$\xi(s)$）和约 $11.3$（$\xi(\mu)$），前者接近合理值，后者仍偏高（可能源于
    仿真器误差或 bin 间相关性估计偏差）。
  \item \textbf{Hartlap 修正}\cite{Hartlap2007}。 当用 $n_{\rm sim}$ 组模拟估计 $p$ 维
    协方差矩阵时，精度矩阵的无偏估计需乘以
    $f_{\rm Hartlap}=(n_{\rm sim}-p-2)/(n_{\rm sim}-1)$，若模拟数量接近 bin 数量该因子
    将显著偏离 $1$。
\end{enumerate}
以上两点意味着本文的 FoM 数值不具有绝对意义，不代表 CSST 的实际约束能力。由于协方差矩阵
对所有核函数完全相同，核函数间 FoM 的相对排序仍然有效——但需注意：当
$\chi^{2}/{\rm dof}\ll 1$（似然面极平坦）时，FoM 差异可能部分反映仿真器在先验边界附近的
外推行为差异，而非纯粹的核函数预测能力差异。本文的 MAE 比较（\autoref{tab:loo}、
\autoref{tab:wilcoxon-vs-m52}）不受此问题影响，是核函数性能排序最可靠的依据。在实际
CSST 数据分析中，协方差矩阵应从与观测体积匹配的模拟中估计，并施加 Hartlap 修正。

此外，$\sigma_8$ 的偏移方向在不同核函数和约束模式间不一致：使用 $\xi(s)$ 的约束 $\sigma_8$
偏低约 $0.5$--$0.7\sigma$（如 RBF $\xi(s)$: $0.775$，偏离 fiducial $0.810$ 约 $-0.47\sigma$），
而 RBF 的 $\xi(\mu)$ 单独约束偏高（$0.856$，$+0.25\sigma$），Mat\'ern 的 $\xi(\mu)$ 则
大幅偏低（$0.644$，$-1.56\sigma$）。这种方向不一致性进一步表明，当前约束更多反映了
协方差矩阵高估和仿真器外推误差的系统效应，而非真实的参数约束。所有偏移均在 $2\sigma$
以内，在统计上不显著。

\section{联合约束与信息冗余分析}\label{sec:joint-info}

RBF 核的联合约束 FoM（$49.3$）与 $\xi(s)$ 单独完全一致，这并非代码错误，而是因为 RBF 核的
$\xi(\mu)$ 仿真器完全丧失了参数敏感性。诊断分析表明：RBF 核的 $\xi(\mu)$ 模型中
WhiteKernel 噪声项占总方差的 $99.93\%$（${\rm noise\_level}=0.993$），信号核方差仅为
$0.026^{2}=0.0007$，且 RBF length\_scale 被推到极小值（$0.011$），无法建立有效的空间相关性。
结果是仿真器对任意参数输入均给出恒定的预测值（等效于训练集加权平均），$\xi(\mu)$ 的似然
函数在参数空间中为常数，对 MCMC 后验无任何约束贡献。

相比之下，Mat\'ern 核的 $\xi(\mu)$ 仿真器保留了微弱但非零的参数敏感性（$\chi^{2}$ 随
$\Omega_c$ 变化范围约 $2.3$），使其联合 FoM（$58.9$）相比 $\xi(s)$ 单独（$47.2$）提升
$25\%$。这揭示了核函数选择不仅影响预测精度（MAE），更可能决定仿真器能否``看到''参数
依赖性本身。

需要直面的问题是：在当前数据条件下（逐样本均值归一化、仅 129 组训练样本），$\xi(\mu)$ 的
GP 仿真器是否根本不可行？答案是有条件的：Mat\'ern 核证明了弱参数依赖性仍可被学习
（FoM 从 $47.2$ 提升至 $58.9$），但 RBF 核的完全失效表明这极度依赖核函数选择。改善方向
包括：（1）取消逐样本均值归一化以保留振幅信息；（2）采用 ARD 核让 GP 学习参数敏感度差异；
（3）增加训练样本数量。在这些改进实现之前，基于当前 $\xi(\mu)$ 仿真器的 MCMC 约束结论
应谨慎解读。

对原始 129 组 N 体模拟数据（非仿真器预测）的 PCA 分析揭示了信息结构：$\xi(s)$ 的参数
依赖性高度集中在第一主成分（解释 $93.3\%$ 方差），而 $\xi(\mu)$ 更分散（PC1 仅解释
$13.4\%$ 方差）；两者 PC1 的相关系数约 $-0.47$，表明 $\xi(\mu)$ 包含与 $\xi(s)$ 部分
互补的信息。$\xi(\mu)$ 在联合约束中贡献有限的主要原因是其信噪比低
（仿真器 MAE/统计误差 $=15.8\%$），加之逐样本均值归一化削弱了振幅信息，而非两者信息
完全冗余。

\section{与已有仿真器工作的比较}\label{sec:compare-prior}

Cosmic Emulator\cite{Heitmann2009}报告物质功率谱精度约 $1\%$，
EuclidEmulator2\cite{Knabenhans2019}在 $z=0$ 的精度优于 $1\%$。本文 $\xi(s)$ 仿真器 MAE
约 $0.5\%$（信号幅度占比），达到或超过同等精度水平。$\xi(\mu)$ 精度 $1.5\%$ 略低，但
考虑其信号动态范围极小（$0.87$--$1.28$），这是合理的。需指出：不同仿真器的统计量、
参数维度和训练集规模各异，精度数字不宜直接比较。本文贡献在于为给定仿真器框架提供核函数
选择的定量依据。

\section{计算可扩展性}\label{sec:scalability}

GP 回归的核矩阵求逆复杂度为 $O(n^{3})$。本文 129 个训练样本完全可行（单次训练约 $2$ 秒）。
当训练集扩展到 1000+ 样本时，可考虑：（1）稀疏 GP（FITC/SVGP），通过诱导点将复杂度降至
$O(nm^{2})$；（2）随机傅里叶特征（RFF），将核近似为有限维内积；（3）神经网络仿真器
（MLP/CNN），在大数据量下可能更具优势。本文核函数比较结论（Mat\'ern 优于 RBF）在上述
近似方法中的适用性值得进一步研究。

\section{局限性}\label{sec:limitations}

\begin{enumerate}
  \item \textbf{参数约束简化。} 仅拟合 $\Omega_c$ 和 $\sigma_8$ 两个参数，余下 7 个固定在
    fiducial 值；实际分析边缘化所有参数后 FoM 将显著降低。本文 FoM 仅作核函数相对比较基准。
  \item \textbf{重复性。} 3-seed 验证显示 $\xi(s)$ 种子间标准差 ${<}10^{-6}$（完全稳定），
    $\xi(\mu)$ 约 $0.0001$--$0.0002$（波动极小）。20 次 Optuna 试验对 $\xi(s)$ 已充分收敛；
    对 $\xi(\mu)$ 而言可通过增加 ${\rm n\_trials}$ 进一步减小波动。
  \item \textbf{协方差矩阵。} 体积不匹配与 Hartlap 修正问题（见 \autoref{sec:omega-c-bias}）
    影响约束结果的绝对可信度，但不影响核函数间的相对比较结论。
  \item \textbf{观测值外推问题。} MCMC 中使用的``观测值''来自 $600~\mathrm{Mpc}/h$ 模拟箱，
    与训练数据的 $1000~\mathrm{Mpc}/h$ 箱不同，导致约 $51.5\%$ 的 $\xi(\mu)$ bin 落在
    训练数据范围之外，仿真器被迫外推。这是 MCMC 约束结果（特别是 $\Omega_c$ 偏移）的
    重要系统误差来源，但不影响核函数间的相对比较。
  \item \textbf{超参数优化的间接信息泄露。} Optuna 在测试集上优化超参数，而 LOO-CV 使用
    该超参数评估，存在间接泄露（详见 \autoref{sec:hpo} 讨论）。
  \item \textbf{MCMC 先验超出训练边界。} MCMC 先验 $\Omega_c\in[0.2,0.4]$，但训练数据
    $\Omega_c$ 仅覆盖 $[0.187,0.352]$，约 $11$--$14\%$ 的后验样本落在训练范围之外，
    仿真器在参数空间中也存在外推。
  \item \textbf{各向同性 length\_scale。} 本文 GP 使用标量 $\ell$（各向同性核），假设
    9 个参数对统计量的影响``等距''。采用 ARD（Automatic Relevance Determination，每维
    独立 length\_scale）可能显著改善性能，是未来重要改进方向。
  \item \textbf{length\_scale 初始值敏感性。} 本文 Optuna 超参数搜索仅优化
    ${\rm constant\_value}$、${\rm noise\_level}$ 和 $\alpha$ 三个参数，length\_scale
    的初始值固定为 $1.0$，其最终取值依赖 sklearn 内层 L-BFGS-B 优化器的调整。不同核函数
    对 length\_scale 初始值的敏感度可能存在差异——例如 Mat\'ern$(\nu=1/2)$ 的似然面可能
    比 RBF 更加多峰，固定初始值可能导致优化器陷入不同局部极值。本文未对此进行系统性消融
    实验（如多起点随机重启），因此不能完全排除 length\_scale 初始值选择对核函数排名的
    潜在影响。未来工作可将 length\_scale 初始值纳入 Optuna 外层搜索空间，或采用多起点
    策略以提高优化的鲁棒性。
  \item \textbf{未考虑实际系统效应。} 选择效应、光纤碰撞、测光红移误差等均未纳入分析。
\end{enumerate}
